\scanpage{94}
une composante $D_\alpha=\overline{\{s_\alpha\}}$ de ce diviseur, n'est autre que sa \uline{multiplicité} en $D_\alpha$.

Ceci rappelé, on se demande donc de \unsure{ce|en} qui se passe pour les $s\in\bigcup D_\alpha$, \deleted{\unsure{On|(les}} \inserted{et} \deleted{voit que} \inserted{\unsure{regarder.|regardons}} ceux qui sont $\in D_\alpha$, $\alpha$ \unsure{pris|fixés|fixé}, i.e. qui spécialisent un $s_\alpha$ fixé. On peut dire que la situation définit une extension $\mathcal E_{U_{s_\alpha}}$ de $\mathcal E_{k(s_\alpha)}$ par $\pi$ --- déduite, comme \unsure{tantôt|tantot} pour $S$ et $K$, d'une courbe \unsure{relatif|relative} $U_{D_\alpha}$ sur $D_\alpha$ --- ou $\unsure{U_{S_\alpha}|U_{s_\alpha}}$ sur un ouvert régulier convenable de $D_\alpha$, si on préfère --- en passant au pt générique. Cette fois-ci, pour la fibre générique, on n'a plus des \deleted{\unsure{sections|section}} pt. rat. sur $k(s_\alpha)$ de $\unsure{U_{s_\alpha}|U_{S_\alpha}}$ --- on n'a qu'un pt : l'infini $i_\alpha\in X_{s_\alpha}\setminus U_{s_\alpha}$, rat sur $\unsure{k_\alpha\xi^{\cdot}|K_\alpha\xi^{\cdot}}$. Mais on dispose \inserted{maintenant} d'un $L_i\subset\pi$ relatif : cet $i_\alpha$, justement, de façon ess. canonique,
\margin{93}
